\section{State-of-the-Art in Interface Prediction} \label{section:state_of_the_art}

Recent years have seen rapid developments in the modelling of interface structures, driven both by methodological advances and by the growing technological importance of heterostructured materials. Interface prediction now spans a range of computational approaches, from first-principles calculations and high-throughput screening to data-driven and machine learning (ML) methods. This section provides an overview of the current state-of-the-art, establishing the context for the subsequent discussion of recent studies (Section 2.5.1), the positioning of this project\rqss methodology (Section 2.5.2), and anticipated emerging directions (Section 2.5.3).

A key enabler of recent progress has been the development of automated interface generation tools such as ARTEMIS and InterMatch, which systematically enumerate possible lattice matches, terminations, and stacking configurations between two crystal structures. When integrated with rapid energy evaluation, either via density functional theory (DFT) or machine-learned potentials (MLPs), these tools allow the configurational space of interface structures to be sampled more broadly and with greater fidelity \cite{Maji2021-kc, Leitherer2023-ew, Lorentzon2025-si, Gerber2023-yd, Rosen2022-pw, Butler2019-ws, Di_Liberto2022-gz, Chagarov2015-ol, Chen2025-bo, Taylor2023-wh, Ceperley1980-bs}. In parallel, studies have increasingly sought to move beyond idealised, coherent interfaces to include semi-coherent, faceted, and amorphous boundaries, reflecting the complexity of real materials under fabrication conditions \cite{Leung2020-tv, Butler2019-ws, Chen2025-bo}.

Efforts have also been made to improve the prediction of interfacial electronic properties, such as band alignment. Where traditional rules like Anderson\rqss rule fail, corrective models based on physically grounded terms, such as $\Delta E_{\Gamma}$ and $\Delta E_{\mathrm{IF}}$, have been proposed and validated across a range of 2D materials \cite{Maji2021-kc, Low2025-on, Rosen2022-pw, Butler2019-ws, Chen2025-bo, Davies2021-zh, Taylor2023-wh}. Such corrections are essential for capturing charge transfer and hybridisation effects at interfaces, and form an important bridge between local structure and electronic response \cite{Davies2024-ju, Leung2020-tv, Butler2019-ws, Chen2025-bo, Davies2021-zh, Taylor2023-wh}.

This chapter aims to synthesise these various advances, situating the present work within an evolving landscape of interface modelling. Particular attention will be given to scalable workflows that combine structural generation, energetic evaluation, and property prediction, setting the stage for a more autonomous and predictive framework for materials interface discovery.

\subsection{What recent work has been done in this area?}

Recent years have witnessed a pronounced expansion in methodologies for predicting interface structures, driven by the convergence of first-principles modelling, high-throughput screening, and data-driven inference.

This evolution reflects a shift from idealised, geometry-constrained models towards scalable frameworks capable of capturing configurational complexity, long-range strain, and emergent phenomena at interfaces \cite{Maji2021-kc, Leitherer2023-ew, Nykiel2023-jf, Rosen2022-pw, Butler2019-ws, Di_Liberto2022-gz, Chen2025-bo}.

\subsubsection{First-principles structure prediction}

Hybrid methodologies now dominate the landscape of ab initio interface prediction. Notable among them is the \textsc{ARTEMIS} tool developed by Taylor et al., which automates the generation of candidate interface structures by identifying lattice-matched slabs, multiple surface terminations, and lateral shifts across a wide range of Miller planes. This approach enables high-throughput sampling of plausible interfaces for subsequent relaxation using density functional theory (DFT), bridging the gap between idealised registry-based models and defect-tolerant configurations \cite{Leung2020-tv, Leitherer2023-ew, Lorentzon2025-si, Nykiel2023-jf, Rosen2022-pw, Butler2019-ws, Di_Liberto2022-gz, Chen2025-bo, Taylor2023-wh}.

Shortly after, Taylor et al.\ introduced the \textsc{RAFFLE}, with Pitfield et al.\ introducing methodology, combining empirical void-filling algorithms with iterative statistical descriptors to model metastable and low-symmetry reconstructions. Their work, particularly on graphene-encapsulated MgO, demonstrates that confinement and bonding at the interface can stabilise novel phases, such as monolayer rocksalt MgO, not observed in bulk form \cite{Hepplestone2005-zl, Hepplestone2006-lo, Butler2019-ws, Chen2025-bo, Taylor2023-wh}.

\subsubsection{Search-based algorithms}

Global optimisation methods have been adapted to the interfacial context. Constrained variants of the Minima-Hopping Method (MHM) allow for controlled exploration of grain boundary reconstructions by varying atomic densities and introducing geometric constraints. These approaches yield more realistic reconstructions than fixed-stoichiometry DFT or classical potentials, especially in systems like silicon where interface states are sensitive to atomic arrangement and misfit accommodation \cite{Anderson1960-zp, Hepplestone2006-lo, Butler2019-ws, Chen2025-bo}.

\subsubsection{Machine learning integration}

Machine learning (ML) is increasingly embedded into the prediction workflow. Gerber et al.\ developed the \textsc{InterMatch} framework, which leverages high-throughput DFT databases to train ML models that guide lattice matching and stacking predictions in 2D systems \cite{Gerber2023-yd, Butler2019-ws, Chen2025-bo}. At the electronic level, Davies et al.\ proposed physically motivated corrections to band alignment, namely $\Delta E_\Gamma$ and $\Delta E_{\mathrm{IF}}$, to extend Anderson\rqss rule and account for hybridisation and interface dipoles in 2D heterostructures \cite{Davies2024-ju, Leung2020-tv, Low2025-on, Butler2019-ws, Chen2025-bo, Davies2021-zh, Taylor2023-wh}.

Moreover, the use of ML surrogates such as MACE and ChgNet has enabled large-scale relaxation of complex interfaces. These interatomic potentials retain near-DFT accuracy while significantly lowering the computational cost, allowing relaxation of thousands of atoms and exploration of non-ideal motifs including semi-coherent, faceted, or disordered junctions \cite{Batatia2022-xz, Schutt2019-sz, Khot2025-hh, Maji2021-kc, Hepplestone2005-zl, Leung2020-tv, Leitherer2023-ew, Hepplestone2006-lo, Lorentzon2025-si, Kao2025-ay, Nykiel2023-jf, Uhrin2025-ba, Owen2024-wp, Rosen2022-pw, Butler2019-ws, Di_Liberto2022-gz, Chen2025-bo, Ceperley1980-bs, Behler2007-nu, Turlo2019-pd, Pitfield2025-ah}.

\subsubsection{Ongoing challenges}

Despite these advances, several limitations persist. Data scarcity for interfacial systems hampers generalisation, and most ML models are trained on bulk-like environments. Existing descriptors may fail to capture the broken symmetry, local dipoles, and charge redistribution unique to interfaces. Furthermore, many workflows truncate the configurational space prematurely or rely on geometric heuristics alone.

There is increasing recognition that hybrid workflows, combining ML surrogates with tools like \textsc{ARTEMIS} and \textsc{RAFFLE}, and validating critical configurations using DFT, offer a scalable route forward \cite{Maji2021-kc, Leitherer2023-ew, Lorentzon2025-si, Nykiel2023-jf, Rosen2022-pw, Butler2019-ws, Di_Liberto2022-gz, Chen2025-bo, Taylor2023-wh}. This multiresolution strategy enables both efficient exploration and physically grounded evaluation of interfacial stability across diverse systems \cite{Khot2025-hh, Leung2020-tv, Butler2019-ws, Chen2025-bo}.

\subsection{How does my proposed approach compare and build upon it?}

Contemporary approaches to interface structure prediction combine symmetry-based lattice matching, high-throughput screening, and density functional theory (DFT) relaxation. Tools such as \textsc{ARTEMIS} enable the systematic generation of interface candidates by identifying lattice commensurabilities and constructing multiple terminations and stacking arrangements from parent surfaces. However, the energetic evaluation of these candidates remains a bottleneck: full relaxation and total energy comparison using DFT is prohibitively expensive for systems exceeding a few hundred atoms \cite{Nykiel2023-jf, Butler2019-ws, Chen2025-bo, Perdew1996-no}.

To mitigate this, recent advances have introduced surrogate strategies that bypass exhaustive DFT sampling. For example, Davies \textit{et al.} proposed semi-empirical corrections (\(\Delta E_\Gamma\), \(\Delta E_{\mathrm{IF}}\)) to Anderson\rqss rule for predicting band alignment in 2D heterostructures, based on descriptors derived from the constituent materials \cite{Davies2024-ju, Butler2019-ws, Davies2021-zh, Taylor2023-wh}. InterMatch, a high-throughput framework proposed by Gerber \textit{et al.}, guides stacking predictions using pre-computed DFT databases and electrostatic heuristics \cite{Maji2021-kc, Gerber2023-yd, Butler2019-ws}. Yet such frameworks often lack spatial resolution or generalisability across chemically diverse or defective systems.

This project builds directly upon these developments, advancing them in three ways:

\begin{enumerate}
   \item \textbf{Surrogate model integration.} Machine-learned interatomic potentials (MLPs), such as MACE, are employed to relax interfacial geometries with near-DFT accuracy but at significantly reduced cost \cite{Batatia2022-xz, Schutt2019-sz, Leitherer2023-ew, Lorentzon2025-si, Nykiel2023-jf, Uhrin2025-ba, Owen2024-wp, Rosen2022-pw, Butler2019-ws, Di_Liberto2022-gz, Ceperley1980-bs, Behler2007-nu, Turlo2019-pd, Pitfield2025-ah}. This enables structural optimisation of larger supercells that accommodate semi-coherent boundaries, strain gradients, or reconstructed motifs inaccessible to traditional DFT-only workflows \cite{Khot2025-hh, Hepplestone2005-zl, Leung2020-tv, Anderson1960-zp, Hepplestone2006-lo, Kao2025-ay, Butler2019-ws, Taylor2023-wh}.
   \item \textbf{Descriptor-led generalisation.} Inspired by the success of band-alignment corrections in TMDC heterostructures, this work seeks to identify whether interfacial stability can be inferred from properties computable on isolated surfaces; such as cleavage ene.g. surface dipole, or projected local density of states near the Fermi level. If validated, such descriptors would permit extrapolation to unseen terminations or chemistries without requiring full interface models \cite{Leung2020-tv, Butler2019-ws, Davies2021-zh, Taylor2023-wh}.
   \item \textbf{Feedback into structure-generation pipelines.} In the longer term, this project aims to integrate ML-derived stability metrics into tools like \textsc{ARTEMIS} or \textsc{RAFFLE}, allowing early pruning of unfavourable configurations before computationally costly evaluation. This hybrid generation--relaxation loop reflects a broader trend towards data-assisted automation in interfacial materials modelling \cite{Leung2020-tv, Butler2019-ws, Davies2021-zh, Taylor2023-wh}.
\end{enumerate}

The proposed framework thus combines multi-resolution structure prediction with predictive modelling rooted in physically meaningful descriptors. Rather than replacing existing DFT-centric approaches, it seeks to augment them; by embedding surrogate models within the structure-generation workflow and using them to navigate large configurational spaces. This enables a shift from retrospective evaluation to prospective design of realistic, energetically favourable interfaces.

\subsection{Emerging trends}

Recent advances in interface prediction have increasingly centred on integrating data-driven approaches with atomistic modelling frameworks. Machine-learned potentials (MLPs) such as MACE and E(3)-equivariant networks like NequIP are now enabling the efficient evaluation of interfacial configurations across large configurational spaces \cite{Batatia2022-xz, Uhrin2025-ba, Owen2024-wp, Rosen2022-pw, Butler2019-ws, Di_Liberto2022-gz, Ceperley1980-bs, Behler2007-nu, Turlo2019-pd, Pitfield2025-ah}. This shift permits relaxation of structures at near-DFT accuracy, allowing researchers to explore interface separations, registry shifts, and metastable motifs that would be prohibitively expensive using standard ab initio methods.

Alongside these developments, tools like ARTEMIS and RAFFLE have emerged to systematically generate candidate interfaces. These platforms sample multiple surface terminations and stackings while accounting for lattice mismatch, enabling more representative modelling of real, faceted, or incoherent boundaries. The capacity to generate and filter hundreds of interfaces also creates opportunities for coupling with high-throughput screening and surrogate models \cite{Leung2020-tv, Leitherer2023-ew, Lorentzon2025-si, Nykiel2023-jf, Butler2019-ws, Taylor2023-wh}.

Emergent trends also include the move toward physically motivated corrections for electronic properties at interfaces. For instance, the failure of Anderson\rqss rule in 2D heterostructures has prompted the introduction of corrective terms such as $\Delta E_\Gamma$ and $\Delta E_\mathrm{IF}$, which better account for hybridisation and interface dipoles. These terms enable more accurate band alignment predictions across families of 2D materials without requiring full-scale DFT recalculation for each pair \cite{Davies2024-ju, Leung2020-tv, Butler2019-ws, Davies2021-zh, Taylor2023-wh}.

Collectively, these approaches suggest a field increasingly focused on scalable, automatable, and physically interpretable modelling of interfaces. The combination of machine learning with established electronic structure theory is opening a path toward predictive frameworks that are sensitive to real-world synthesis constraints and device contexts \cite{Khot2025-hh, Leung2020-tv, Butler2019-ws}.